% 6-glossary.tex starts here
% Requires in the preamble: \usepackage{hyphenat}
% ----------------------------------------------------------
% ACRONYMS (used with \gls/\acr... rendered with \printglossary[type=\acronymtype])
% ----------------------------------------------------------

\newacronym{pan}{PAN}{Prisma–Analog–Nx}

\newacronym{ssr}{SSR}{server-side rendering}
\newacronym{ssg}{SSG}{static site generation}
\newacronym{csr}{CSR}{client-side rendering}
\newacronym{dsn}{DSN}{Data Source Name}

\newacronym{trpc}{tRPC}{TypeScript remote procedure call}

\newacronym[longplural={data transfer objects}, shortplural={DTOs}]
  {dto}{DTO}{data transfer object}

\newacronym[longplural={object–relational mappings}, shortplural={ORMs}]
  {orm}{ORM}{object–relational mapping}


\newacronym{e2e}{E2E}{end-to-end testing}

\newacronym{ci}{CI}{continuous integration}
\newacronym{cicd}{CI/CD}{continuous integration and continuous delivery/deployment}

\newacronym{p50}{p50}{50th percentile (median)}
\newacronym{p95}{p95}{95th percentile}

\newacronym{lcp}{LCP}{Largest Contentful Paint}
\newacronym{cls}{CLS}{Cumulative Layout Shift}
\newacronym{fcp}{FCP}{First Contentful Paint}
\newacronym{tbt}{TBT}{Total Blocking Time}
\newacronym{avif}{AVIF}{image format based on the AV1 codec}

\newacronym{csp}{CSP}{Content-Security-Policy}
\newacronym{hsts}{HSTS}{HTTP Strict-Transport-Security}

\newacronym[longplural={minimum viable products}, shortplural={MVPs}]
  {mvp}{MVP}{minimum viable product}


\newacronym[longplural={user interfaces}, shortplural={UIs}]
  {ui}{UI}{user interface}


\newacronym{mean}{MEAN}{MongoDB–Express–Angular–Node.js}


\newacronym{spa}{SPA}{single-page application}
\newacronym{seo}{SEO}{search engine optimization}
\newacronym[longplural={pull requests}, shortplural={PRs}]
  {pr}{PR}{pull request}


% ----------------------------------------------------------
% TERMS and NAMES (used with \gls rendered with \printglossary[title={Glossary of Terms}])
% ----------------------------------------------------------

\newglossaryentry{web}{
  name={\NoHyphens{web}},
  sort={web},
  description={the ecosystem of \emph{hypertext} resources accessible over the Internet (World Wide Web); in this thesis, “web application” denotes an application accessed from the browser, with \glsxtrshort{ssr}/\glsxtrshort{ssg}/\glsxtrshort{csr} rendering},
}

\newglossaryentry{nx}{
  name={\NoHyphens{Nx}},
  description={build and orchestration system for JavaScript/TypeScript monorepos},
}

\newglossaryentry{prisma}{
  name={\NoHyphens{Prisma}},
  description={\glsxtrshort{orm} library for TypeScript/Node.js},
}

\newglossaryentry{lighthouse}{
  name={\NoHyphens{Lighthouse}},
  description={Google tool for automated performance auditing of \gls{web} pages, which reports \emph{Core Web Vitals} metrics (\glsxtrshort{lcp}, \glsxtrshort{cls}, \glsxtrshort{fcp}, \glsxtrshort{tbt})},
}

\newglossaryentry{analog}{
  name={\NoHyphens{Analog}},
  description={full-stack meta-framework for \gls{angular}, with \glsxtrshort{ssr} support and file-based routing; used for \gls{web} applications},
}

\newglossaryentry{angular}{
  name={\NoHyphens{Angular}},
  description={platform and framework for \gls{web} applications (SPA/\glsxtrshort{ssr})},
}

\newglossaryentry{nitro}{
  name={\NoHyphens{Nitro}},
  description={universal runtime (UnJS) for server routes and server-side rendering (\glsxtrshort{ssr})},
}


\newglossaryentry{zod}{
  name={\NoHyphens{Zod}},
  description={schema-based validation library for TypeScript/JavaScript},
}

\newglossaryentry{cockroachdb}{
  name={\NoHyphens{CockroachDB}},
  description={distributed (NewSQL) database management system, compatible at the protocol level with \gls{postgresql}, with automatic replication and partitioning for resilience and horizontal scaling},
  sort={cockroachdb},
}

\newglossaryentry{postgresql}{
  name={\NoHyphens{PostgreSQL}},
  description={open-source relational database management system, with extensive support for the SQL standard, ACID transactions and extensibility through user-defined types and functions},
  sort={postgresql},
}


\newglossaryentry{vercel}{
  name={\NoHyphens{Vercel}},
  description={cloud platform for building, previewing and deploying \gls{web} applications, with \glsxtrshort{cicd} integration, per-environment variables and support for \glsxtrshort{ssr}, \textit{serverless} functions and \textit{edge functions}},
  sort={vercel},
}

\newglossaryentry{vitest}{
  name={\NoHyphens{Vitest}},
  description={fast testing framework for the Vite ecosystem, aimed at unit and integration tests, with an API largely compatible with Jest and native TypeScript support},
  sort={vitest},
}

\newglossaryentry{playwright}{
  name={\NoHyphens{Playwright}},
  description={library for browser automation and \glsxtrshort{e2e} testing, with support for Chromium, Firefox and WebKit, headless/GUI execution and tools such as the trace viewer},
  sort={playwright},
}

\newglossaryentry{hydration}{
  name={\NoHyphens{hydration}},
  description={the process by which client-side code attaches interactive behavior to the \emph{markup} generated through \glsxtrshort{ssr}; it allows the transition from the initial server-side rendering to \glsxtrshort{csr} interactions without reloading the page},
  sort={hydration},
}

\newglossaryentry{monorepo}{
  name={\NoHyphens{monorepo}},
  description={a code organization strategy in which several applications and libraries coexist in a single repository, with unified orchestration of builds, tests and \glsxtrshort{cicd}; in the context of this thesis it is managed with Nx},
  sort={monorepo},
}

\newglossaryentry{uses}{
  name={\NoHyphens{USES}},
  sort={uses},
  description={a relation between modules defined by Parnas (1972): a module A is in the \textit{USES} relation with module B if the correct functioning of A depends on the presence of a correct version of B. Parnas writes it in capital letters as a formal term and requires this relation to be a loop-free (acyclic) partial ordering, a property that allows functional subsets of the system to be built by “cutting off” the upper levels},
}


\newglossaryentry{typescript}{
  name={\NoHyphens{TypeScript}},
  sort={typescript},
  description={a typed superset of the JavaScript language that introduces static types and is transpiled to JavaScript},
}

\newglossaryentry{kookio}{
  name={\NoHyphens{Kookio}},
  sort={kookio},
  description={full-stack \gls{web} application used as a case study in this thesis for pantry management, daily meal planning and automatic derivation of the shopping list, built on the \gls{pan} architecture},
}


\newglossaryentry{nextjs}{
  name={\NoHyphens{Next.js}},
  sort={nextjs},
  description={React framework oriented towards hybrid rendering (\glsxtrshort{ssr}/\glsxtrshort{ssg}), often associated with JAMstack},
}


\newglossaryentry{vite}{
  name={\NoHyphens{Vite}},
  sort={vite},
  description={build tool and development server for modern \gls{web} applications, based on ECMAScript modules, with very fast startup and \emph{hot module replacement} (HMR); it underpins \gls{vitest} and is used by \gls{analog} for \glsxtrshort{ssr} and builds},
}

\newglossaryentry{pino}{
  name={\NoHyphens{Pino}},
  sort={pino},
  description={performance-oriented logging library for Node.js (JSON logs, low-overhead output)},
}

\newglossaryentry{sentry}{
  name={\NoHyphens{Sentry}},
  sort={sentry},
  description={production error monitoring platform, which aggregates, groups and alerts on the exceptions captured from the client and the server},
}


\newglossaryentry{workflow}{
  name={\NoHyphens{workflow}},
  sort={workflow},
  plural={\NoHyphens{workflows}},
  description={a declarative configuration describing a sequence of automated steps (build, testing, lint, deploy) triggered by events (e.g. push, pull request, schedule). In ecosystems such as GitHub Actions and GitLab \glsxtrshort{ci}, a workflow is made up of jobs and steps, and can use matrices, caches, artifacts, secrets/variables and concurrency rules. In this project, the workflows orchestrate \glsxtrshort{cicd} and align execution with the \gls{nx} targets.},
  see={[see also]{cicd,nx}},
}

\newglossaryentry{compositeaction}{
  name={\NoHyphens{composite action}},
  sort={compositeaction},
  plural={\NoHyphens{composite actions}},
  description={a type of reusable action in GitHub Actions that groups several steps (\texttt{run}/\texttt{uses}) into a single action; it is distributed as a directory with an \texttt{action.yml} file, accepts \textit{inputs}, \textit{outputs} and \textit{secrets} and can be called from a \gls{workflow} through \texttt{uses:}. Useful for eliminating duplication (DRY) in \glsxtrshort{cicd} pipelines; common limitations: they do not run in their own containers and cannot define \textit{services}.},
  see={[see also]{workflow,cicd}},
}


\newglossaryentry{redis}{
  name={\NoHyphens{Redis}},
  sort={redis},
  description={an \emph{in-memory} key–value data store, with advanced structures (lists, sets, hashes, sorted sets, \textit{streams}) and very low latency; frequently used for caching, session storage, queues/messaging (Pub/Sub, Streams), rate-limiting and \textit{leaderboards}. It supports optional persistence (RDB/AOF), replication and \textit{cluster} mode for partitioning and, as a rule, complements a primary SQL database in microservice and \textit{serverless} architectures.},
}


\newglossaryentry{metaframework}{
  name={\NoHyphens{meta-framework}},
  sort={metaframework},
  description={a framework that extends a base framework with full-stack capabilities (e.g. file-based routing, \glsxtrshort{ssr}/\glsxtrshort{ssg}, integration with \glsxtrshort{trpc}). In this thesis, \gls{analog} extends \gls{angular} on top of the \gls{nitro} runtime.},
  see={[see also]{analog,angular,nitro,ssr,ssg,trpc}},
}


\newglossaryentry{randare}{
  name={\NoHyphens{rendering}},
  sort={rendering},
  description={the process by which an application generates HTML (and interface updates) from code and data. On the modern web it can take place on the server (\glsxtrshort{ssr}), at build time through \glsxtrshort{ssg}/prerendering, or in the browser (\glsxtrshort{csr}); as a rule, it is followed by \gls{hydration} for interactivity. In \gls{analog}, the rendering strategy can be chosen per route.},
  see={[see also]{ssr,ssg,csr,hydration,analog}},
}


\newglossaryentry{logging}{
  name={\NoHyphens{logging}},
  sort={logging},
  description={collecting and recording application events (structured messages by level, correlation through identifiers, export to aggregation systems) for debugging and observability},
}


% darwin-arm64 / rhel-openssl-3.0.x glossary entries removed: Prisma 7 is
% Rust-free (no platform-specific query-engine binaries / binaryTargets), so
% the terms no longer appear anywhere in the body. \glsaddall would otherwise
% force-print them as orphaned entries describing a mechanism the project no
% longer uses.
\newglossaryentry{prismaclient}{
  name={\NoHyphens{Prisma Client}},
  sort={prismaclient},
  description={\texttt{TypeScript} database client, generated automatically by \gls{prisma} from the \texttt{schema.prisma} file; it provides a typed API for SQL operations (CRUD, filters, relations, transactions) on the defined models. Starting with \gls{prisma} 7, the client runs entirely in JavaScript (without a native Rust engine): the connection goes through a \textit{driver adapter} (\texttt{@prisma/adapter-pg}), and regenerating with \texttt{npx prisma generate} produces only typed code, independent of the execution platform.},
  see={[see also]{prisma,orm,cockroachdb,postgresql}},
}

% 6-glossary.tex ends here
